\section{Two-endpoint interpolation and exact prime deletion}
\label{sec:endpoint-filter-case}

These are two distinct ways to design an approximation before taking
its primitive form. Endpoint interpolation forces zeros in an integral
source. Prime deletion instead removes selected coefficient denominators
from a generating function. Both operations are exact, but neither
alone controls the normalized error. We give enough information to
construct the endpoint pairs without recovering a large historical
linear system. The records are \recordid{H01}--\recordid{H03} and
\recordid{A04}--\recordid{A05}.

\subsection{A four-by-five system for all rational Hermite sources}
Write $\mathcal F(x)=\sum_{j\geq1}\chi_4(j)x^j/j^4$ and $E=x\,d/dx$,
so $E^4\mathcal F=x/(1+x^2)$ and $\mathcal F(1)=\theta$.
For a source $f$ put
\[
 \mathcal I(f)=\frac16\int_0^1f(x)(-\log x)^3\dd x.
\]
Fix $N\geq1$, positive $a,b$ with $a+b=2N$, and a quadratic
$D\in\Z[x]$ with $D(0)=1$ and no zero on $[0,1]$.
Seek $A\in\Q$ and $R\in\Q[x]$, $\deg R\leq2N$, $R(0)=0$,
such that, with $G=R/D^N$,
\begin{equation}\label{eq:hermite-source}
 f=\frac{A}{1+x^2}+\frac{E^4G}{x}
       =O(x^a)\text{ at }0,\qquad f=O((x-1)^b)\text{ at }1.
\end{equation}
Four integrations by parts, or the monomial identity
$\mathcal I(E^4x^j/x)=1$ for $j\geq1$, give
\begin{equation}\label{eq:hermite-value}
 \mathcal I(f)=A\theta+G(1).
\end{equation}
For the integration argument the rational $G$ is regular on the
integration interval and vanishes at zero; all logarithmic boundary
terms at zero vanish, and the surviving endpoint term is $G(1)$.
Zeros of $f$ may make its integral small, but its sign and normalization
still require separate analysis.

Define
\[
 R_{\rm low}(x)=[-D(x)^N\mathcal F(x)]_{\leq a},\qquad
 L(z)=\prod_{j=a+1}^{2N}(z-j),\qquad D(x)^N=\sum_\ell d_\ell x^\ell,
\]
where brackets mean Taylor truncation. Let
\[
 T(y)=\frac1{2\cosh y},\qquad
 V_b(y)=-\sum_{m=4}^{b+3}\frac{T^{(m-4)}(0)y^m}{m!}.
\]
All these finite coefficients are rational. For $0\leq r,k\leq3$ set
\begin{align}
 K_{rk}&=\sum_\ell d_\ell
          \left.\frac1{k!}\frac{\dd^k}{\dd z^k}(z^rL(z))\right|_{z=\ell},
 \label{eq:hermite-K}\\
 v_r&=\left.\partial_y^rL(\partial_y)
           \bigl(R_{\rm low}(e^y)-D(e^y)^NV_b(y)\bigr)\right|_{y=0}.
 \label{eq:hermite-v}
\end{align}

\begin{proposition}[Exact endpoint parametrization; H01]\label{prop:hermite-map}
The source space~\eqref{eq:hermite-source} is naturally identified with
\[
 Kg=Av,\qquad g_k=\left.\partial_y^kG(e^y)\right|_{y=0},\quad0\leq k\leq3.
\]
Its endpoint form is $A\theta+g_0$. This includes singular $K$ and
zero-target solutions; no determinant is divided out.
\end{proposition}
\begin{proof}
The condition at zero is exactly
$R=AR_{\rm low}+\sum_{j=a+1}^{2N}r_jx^j$.
Put $H(y)=G(e^y)$. Multiplying the source by $x=e^y$ changes it to
$AT(y)+H''''(y)$, so the condition at one is
\[
 H(y)=\sum_{k=0}^3g_ky^k/k!+AV_b(y)+O(y^{b+4}).
\]
The condition at zero is equivalently
\[
 L(\partial_y)\bigl(D(e^y)^NH(y)-AR_{\rm low}(e^y)\bigr)=0.
\]
Apply $\partial_y^r$, $r=0,1,2,3$, and evaluate at zero. The identity
$P(\partial_y)(e^{\ell y}y^k/k!)|_0=P^{(k)}(\ell)/k!$
gives exactly $Kg=Av$.

Conversely, use the prescribed Taylor polynomial for $H$ as initial
data for the displayed order-$b$ differential equation. Its leading
coefficient at zero is $D(1)^N\ne0$. The four equations say that the
residual of that polynomial vanishes to order four, so the differential
recurrence makes the solution agree through degree $b+3$.
This also handles $b<4$; the same recurrence forces the additional
specified derivatives. The equation implies
$D(e^y)^NH-AR_{\rm low}(e^y)=\sum_{j=a+1}^{2N}r_je^{jy}$.
Its distinct integer exponents and rational initial data give rational
$r_j$ by a Vandermonde system. Thus it reconstructs the required $R$.
\end{proof}

\subsection{Primitive endpoint pairs and the actual error}
Every entry of $K$ is integral, since a divided derivative of an integer
polynomial at an integer is integral. Let $d_a$ be the lcm of the odd
positive integers at most $a$ and $C=2d_a^4$. Then $Cv$ is integral.
For $R_{\rm low}$ this follows from its coefficient denominators;
for $V_b$ it follows from $T^{(j)}(0)$ being half-integers and from
Leibniz's rule, whose evaluated factorial ratios are integers.
Cramer's identity supplies the integer endpoint pair
\begin{equation}\label{eq:hermite-pair}
 (A,B)=\bigl(C\det K,\ \det K[\text{column }0\leftarrow Cv]\bigr).
\end{equation}
If this pair is zero it is not a witness. If the augmented matrix has
rank less than four, compute its full rational kernel instead of using
vanishing maximal minors. For every nonzero pair, divide its complete
gcd before any numerical comparison.

There is an exact formula for the signed error, not merely its source.
Set $S_a(x)=R_{\rm low}(x)+D(x)^N\mathcal F(x)$ and
$\ell_r=\partial_y^rL(\partial_y)S_a(e^y)|_0$.
Since $\partial_y^4\mathcal F(e^y)=T(y)$, the column relation is
$v=\ell-K\eta$, where $\eta_k=\partial_y^k\mathcal F(e^y)|_0$ and
$\eta_0=\theta$. Multilinearity therefore gives
\begin{equation}\label{eq:hermite-error}
 \det K\,\theta+\det K[0\leftarrow v]=\det K[0\leftarrow\ell].
\end{equation}
All lower-period endpoint derivatives disappear through repeated columns.
The identity is valid even when $K$ is singular.

A useful factorial normalization changes neither primitive pair.
With $d_b=\operatorname{lcm}(1,\ldots,b)$,
$d_b^kK_{rk}/b!$ is integral: expand
$L(z)/b!=\binom{z-a-1}{b}$ at an integer by Vandermonde's identity,
and use
$\binom{t}{j}=(-1)^{j-1}(t/j)\prod_{s=1}^{j-1}(1-t/s)$.
The coefficient of $t^k$ has denominator dividing $d_b^k$.
Replacing $K$ by $K\operatorname{diag}(1,d_b,d_b^2,d_b^3)/b!$
and $v$ by $v/b!$ multiplies both determinants by
$d_b^6/(b!)^4$. It improves a representation, not the primitive pair.

For $D=1+x^2$, $a=b=N=n$, the scalar remainder $S_n$ is positive
on $x>0$. To see this, use
$\mathcal F(x)=x\int_0^1(-\log u)^3/(1+x^2u^2)\dd u/6$.
Put $m=\lceil n/2\rceil$, $y=x^2$, $t=u^2$. The remainder of
$(1+y)^n/(1+ty)$ after degree $m-1$ in $y$ is
\[
 \frac{y^m}{1+ty}\left(c_m(t)+\sum_{k=m+1}^n\binom nk y^{k-m}\right),
 \qquad c_m(t)=\sum_{k=0}^m\binom nk(-t)^{m-k}.
\]
In reverse order the alternating terms defining $c_m$ decrease in
magnitude on $0\leq t<1$, giving strict positivity. Integrate the
remainder to obtain $S_n>0$. The determinant in~\eqref{eq:hermite-error}
contains high derivatives of $S_n$; its positivity does not follow.
The unresolved questions for a selected schedule are nonzero endpoint
error and the comparison of that error with the reduced pair's size.
Equations~\eqref{eq:hermite-K}--\eqref{eq:hermite-pair} give a fixed-size
exact procedure for investigating them.

\subsection{Prime deletion changes the generating function}
Return to $\mathcal F(t)=\sum\chi_4(j)t^j/j^4$. For an odd prime $p$,
\[
 \mathcal F(t)-\frac{\chi_4(p)}{p^4}\mathcal F(t^p)
       =\sum_{p\nmid j}\frac{\chi_4(j)t^j}{j^4}.
\]
For a finite set $S$ of odd primes put $P=\prod_{p\in S}p$ and
\begin{equation}\label{eq:deleted-germ}
 G_S(t)=\sum_{(j,P)=1}\frac{\chi_4(j)t^j}{j^4},\qquad
 \alpha_S=\prod_{p\in S}(1-\chi_4(p)/p^4),\qquad G_S(1)=\alpha_S\theta.
\end{equation}
Build the bordered pencil~\eqref{eq:common-J} from the coefficients
of $G_S$, and evaluate its variable at $\alpha_S\theta$.
For each $p\in S$ every coefficient entry is $p$-integral, whereas
$v_p(\alpha_S)\geq-4$. Therefore the actual primitive multiplier of
this endpoint pencil satisfies $v_p(\kappa)\leq4$, independently of
matrix size. The endpoint factor is essential to this conclusion.

For matrix order $n$, delete every odd prime $y<p\leq2n$.
The surviving prefix through degree $2n$ is cleared by
\[
 D_{y,2n}=\prod_{\substack{p\leq\min(y,2n)\\p\ {
m odd}}}
              p^{4\lfloor\log_p(2n)\rfloor}.
\]
Scaling the bordered matrix gives a sufficient overall multiplier
$P^4D_{y,2n}^{n+1}$, whence
\[
 \log^+\kappa\leq4\log P+4(n+1)\pi(y)\log(2n)=o(n^2)
\]
if $\pi(y_n)\log n=o(n)$; $y_n=\sqrt n$ suffices even using
$\pi(y)\leq y$ and $\log P\leq2n\log(2n)$.
This is a genuine arithmetic result with no matching analytic theorem.

The analytic cost is concrete. In the divisor expansion
$G_S=\sum_{d\mid P}\mu(d)\chi_4(d)d^{-4}\mathcal F(t^d)$,
apply $E^4$. At a primitive $4P$th root of unity only $d=P$
can have a pole, since $E^4\mathcal F(t^d)=d^4t^d/(1+t^{2d})$.
Its coefficient is nonzero. Thus growing prime sets create genuine
new singularities. Neither the analytic bound nor the nonvanishing
subsequence for a fixed germ can simply be reused for this varying family.

\subsection{A sparse family whose tiny determinant factor cancels}
There is a useful explicit test of the last warning. Let $M=3^k$,
$n=M-1$, and delete every odd prime at most $B\geq2n$ except 3.
The coefficient prefix is $r(t)+ct^M$, where
\[
 r(t)=\sum_{0\leq a<k}\frac{(-1)^at^{3^a}}{3^{4a}},\qquad
 c=\frac{(-1)^k}{M^4},\qquad \deg r=M/3<M/2.
\]
The polynomials $q(t)=r(t)/t-ct^{M-1}$ and $p(t)=r(t)^2/t$ obey
$q(t)(r(t)+ct^M)-p(t)=-c^2t^{2M-1}$, so they give a diagonal
Pad\'e pair for this prefix. More precisely its bordered determinant is
\begin{equation}\label{eq:sparse-collapse}
 J_{n,S}(X)=(-1)^{n(n-1)/2+1}c^{M-2}
                 \bigl((r(1)-c)X-r(1)^2\bigr).
\end{equation}
Here is a verification of the scalar, which matters for normalization.
Reverse the columns of $T$ and multiply by $(I-\mathrm{shift}_{\rm down})^{-1}$.
The result is $cI+N-\mathbf1v^T$, with
$N_{ij}=r_{M+i-j}$, $v_j=r_{M-j}$, $N^2=0$ and $v^TN=0$.
The determinant lemma gives
$\det T=\delta c^{n-1}(c-r(1))$, $\delta=(-1)^{n(n-1)/2}$.
Also $\operatorname{adj}(T)_{11}=\delta c^{n-1}$ and
$T\mathbf1=ce_1-b$. Using these identities in
$X\det T+b^T\operatorname{adj}(T)b$ proves~\eqref{eq:sparse-collapse}
without assuming $T$ invertible.

The entire factor $c^{M-2}$ disappears after primitive normalization.
The remaining form is
$(r(1)-c)\alpha_S\theta-r(1)^2$. At the natural cutoff $B=2n$
its useful real and primitive estimates remain unresolved. Tiny raw
determinants alone therefore do not settle this prescribed sparse family.
An additional large-cutoff obstruction in the catalog concerns a
different scale; it does not close the natural cutoff.
